\section{Case studies}
\label{sec:cases}

The case studies follow the same pattern. The control goals are identified; they are turned
into a cost; the system is modelled with states, inputs and disturbances; the limits of the
system and the logic of the application become constraints; and the resulting problem is solved
at every step on a receding horizon. All of them are economic MPC problems in the sense of
\cref{sec:economic}, and all use the constructions of \cref{sec:toolbox}. For each of them the
course website has an interactive laboratory that runs the controller in the browser.

\subsection{The charging service area}
\label{sec:servicearea}

\paragraph{Problem.} A service area for the fast charging of electric vehicles draws power from
the grid at a point of connection (POC). Vehicles arrive wanting to recharge as fast as
possible, so even a few of them produce a large power flow, which means a large connection and
a heavy bill under a capacity-based tariff. An energy storage system (ESS) is installed to keep
the power at the point of connection small. On site there may also be a renewable generator.

\paragraph{Variables and model.} With the \emph{consumer sign convention} --- power consumed is
positive --- the quantities at time $i$ of the window opened at $k$ are: the charging demand of
the vehicles $p^{pev}_{\ik{i}}$ and the renewable generation $p^{res}_{\ik{i}}$, both known over
the window (disturbances with a forecast); the storage power $u_{\ik{i}}$, positive when
charging, which is the control; the energy stored $x_{\ik{i}}$, which is the state; and the power
at the point of connection $p_{\ik{i}}$, which is the controlled quantity. The balance at the
point of connection and the dynamics of the storage are
\begin{equation}
  p_{\ik{i}} = p^{pev}_{\ik{i}} + u_{\ik{i}} - p^{res}_{\ik{i}}, \qquad
  x_{\ik{i+1}} = x_{\ik{i}} + \tfrac{T}{3600}\, u_{\ik{i}}, \qquad x_{\ik{0}} = x_k .
\end{equation}

\paragraph{Problem solved at every step.}
\begin{equation}
\begin{aligned}
  \min_{u} \quad & \sum_{i=0}^{N-1}\Big[\alpha\, p_{\ik{i}}^2 + \beta\,\big(x_{\ik{i}} - x^{ref}\big)^2\Big]
                   + \beta_f\,\big(x_{\ik{N}} - x^{ref}\big)^2 \\
  \st \quad & \text{balance, dynamics, } x_{\ik{0}} = x_k, \\
            & u^{\min} \le u_{\ik{i}} \le u^{\max},\quad x^{\min} \le x_{\ik{i}} \le x^{\max},\quad
              p^{\min} \le p_{\ik{i}} \le p^{\max}.
\end{aligned}
\end{equation}
The weight $\alpha$ penalizes the power at the point of connection; $\beta$, deliberately
small, keeps the storage near a reference level so that it is ready for the next event; the
terminal weight $\beta_f$ is taken from the Riccati equation of the unconstrained problem, as
\cref{sec:optimality} prescribes. Since $p$ is an affine function of $u$, the limits on $p$ are
bounds on the control, and the problem is a quadratic program.

\paragraph{Losses and indicators.} A real storage has a charging and a discharging efficiency,
and the two cases obey different equations. The storage power is split into a charging and a
discharging power, $p^{ess} = p^{ch} - p^{dis}$, both non-negative, with the indicator variables
$\delta^{ch}, \delta^{dis}$, the constraints $p^{ch} \le p^{ch,\max}\delta^{ch}$, $p^{dis} \le
p^{dis,\max}\delta^{dis}$, $\delta^{ch} + \delta^{dis} \le 1$, and the dynamics with the
efficiencies, as in \cref{sec:toolbox}. The problem becomes a mixed-integer quadratic program.
We also distinguished the \emph{control model}, which lives inside the MPC problem and is the
controller's belief about the plant, from the \emph{simulation model}, which replaces the real
system when the controller is tested.

\paragraph{Code.} The implementation is split in two files on purpose: \texttt{main.jl} is the
application, which holds the parameters, builds the profiles and runs the receding-horizon loop
--- measure, call the optimizer, apply the first control, simulate one step, repeat;
\texttt{MPC\_iteration.jl} is the controller, which builds the optimization problem in the
modelling language JuMP and hands it to a solver. Changing solver is one line, and the model
does not change: this separation between how a problem is written and how it is solved is the
reason for using a modelling language at all.

\begin{labbox}
The Service Area Lab runs the controller over a full day against the same area without storage;
the peak at the point of connection halves. The lab also shows the effect of the horizon, of the
forecast, and of the terminal weight. The Julia code of the lesson is on the website, with the
formulation it implements. Among real installations, battery-buffered fast chargers sold for
sites with a weak grid connection are this problem turned into a product.
\end{labbox}

\subsection{Building energy management}
\label{sec:building}

\paragraph{Problem.} A building has photovoltaic panels, non-plannable loads, plannable loads,
an energy storage system and a charging station, all behind one meter at the point of delivery
(\cref{fig:building}). The first goal is to minimize the electricity bill under a time-of-use
tariff, with one price $c^{in}$ for the energy consumed and a lower price $c^{out}$ for the
energy injected into the grid.

\begin{figure}[t]
\centering
\begin{tikzpicture}[x=1cm,y=1cm,>=Stealth,
  box/.style={draw,rounded corners=2pt,minimum width=1.9cm,minimum height=0.8cm,font=\small,fill=white}]
  \draw[rounded corners=4pt,fill=soft] (-0.4,-1.1) rectangle (8.6,3.6);
  \node[font=\small\bfseries,anchor=north west] at (-0.3,3.55) {building};
  \node[box,fill=sun!20] (pv) at (0.9,2.7) {PV};
  \node[box] (npl) at (3.9,2.7) {non-plannable loads};
  \node[box] (pl) at (7.0,2.7) {plannable loads};
  \node[box,fill=ess!20] (ess) at (0.9,0.5) {storage};
  \node[box] (cs) at (3.9,0.5) {charging station};
  \node[box,fill=mpc!15] (mpc) at (7.0,0.5) {MPC};
  \node[box,minimum width=1.1cm] (meter) at (10.3,1.6) {meter};
  \node[draw,cloud,cloud puffs=11,minimum width=2.1cm,minimum height=1.3cm,font=\small] (grid) at (12.9,1.6) {grid};
  \draw[thick] (8.6,1.6) -- (meter);
  \draw[thick,<->] (meter) -- (grid) node[midway,above,font=\small] {$P$};
  \draw[dashed,mpc,->] (mpc.south) -- ++(0,-0.4) -| (ess.south) node[pos=0.75,left,font=\scriptsize,mpc] {$P^{ess}$};
  \draw[dashed,mpc,->] (mpc) -- (cs) node[midway,above,font=\scriptsize,mpc] {$P^{pev}$};
  \draw[dashed,mpc,->] (mpc) -- (pl) node[midway,right,font=\scriptsize,mpc] {$s_{j,\ell}$};
\end{tikzpicture}
\caption{The building of the second case study. The controller acts on the storage, on the
charging station and on the start times of the plannable loads; generation and non-plannable
loads are forecasts; the bill is paid on the power at the meter.}
\label{fig:building}
\end{figure}

\paragraph{Plannable loads.} Each plannable load $\ell$ comes with the preferences of the user,
the earliest start $S_\ell$, the deadline $F_\ell$, the duration $N_\ell$, and with its power
profile $P_\ell$. Load shifting is the decision of when to start it, written with the one-hot
variables $s_{\ik{j},\ell}$ of \cref{sec:toolbox}: the load starts exactly once,
$\sum_{j=S_\ell}^{F_\ell - N_\ell} s_{\ik{j},\ell} = 1$, its power is the shifted profile
$P^{pl}_{\ik{i},\ell} = \sum_j s_{\ik{j},\ell}\,P_\ell(i-j+1)$, and the total plannable power is
$P^{pl}_{\ik{i}} = \sum_\ell P^{pl}_{\ik{i},\ell}$. Note the two time indices: $S_\ell$ and
$F_\ell$ are absolute times, $i|k$ is relative to the window.

\paragraph{The bill.} The power at the meter, $P_{\ik{i}} = P^{npl} + P^{pl} + P^{pev} + P^{ess}
- P^{pv}$, is a cost when positive and a revenue when negative. It is split as $P = P^{in} -
P^{out}$ with $P^{in}, P^{out} \ge 0$, the two forbidden from being positive together by the
indicators $\delta^{in}, \delta^{out}$ and the constraints $P^{in} \le P^{\max}\delta^{in}$,
$P^{out} \le \abs{P^{\min}}\delta^{out}$, $\delta^{in} + \delta^{out} \le 1$. The bill over the
window is then linear:
\begin{equation}
  J^{cost} = \sum_{i=0}^{N-1} \tfrac{T}{3600}\Big(c^{in}_{\ik{i}} P^{in}_{\ik{i}} - c^{out}_{\ik{i}} P^{out}_{\ik{i}}\Big).
\end{equation}

\paragraph{Battery ageing.} A storage system has a price $COST^{ess}$ and a finite number of
full cycles. The energy it will move over its life is $E^{life} = 2\,x^{ess,\max}\cdot
\mathit{cycles}$, and the ratio $c^{dep} = COST^{ess}/E^{life}$ is what one kWh through the
battery costs in terms of life consumed --- a few cents. Both charging and discharging consume
life, so the term is
\begin{equation}
  J^{depr} = \sum_{i=0}^{N-1} c^{dep}\,\big(P^{ch}_{\ik{i}} + P^{dis}_{\ik{i}}\big)\,\tfrac{T}{3600},
\end{equation}
which is $c^{dep}\abs{P^{ess}}$ written without an absolute value, thanks to the split into two
non-negative powers. It acts as a friction: the storage is used only when the benefit exceeds
the loss of life.

\paragraph{The full objective.} With a regulation term on the stored energy and the terminal
value of \cref{sec:economic},
\begin{equation}
  J = J^{cost} + \alpha_x \sum_{i}\big(x^{ess}_{\ik{i}} - x^{ref}\big)^2 + \alpha_{dep}\,J^{depr}
      - c^{avg}\,x^{ess}_{\ik{N}} ,
\end{equation}
subject to the balance at the meter, the dynamics of the storage with its efficiencies, the
measured initial state, the one-hot constraints of the loads, the indicator constraints and the
bounds on all variables. Without the last term the battery ends every day empty, because selling
its energy is a profit inside the window; with it, the energy left is valued at the average
purchase price and the controller keeps it. The weights are chosen as in \cref{sec:economic},
and they can be normalized to sum to one.

\paragraph{What was left out.} The model is electrical only. A real battery management system
also tracks the state of charge in percent, the state of health, which cell balancing protects,
and the state of power, the power the battery can deliver in the near future. The thermal part
of the building --- wall and room temperatures driven by the weather and by heating and
cooling, humidity, air quality, ventilation --- would add comfort terms to the objective and
constraints of its own. The building information model is the source from which the parameters
of all these models can be updated as the building changes.

\begin{labbox}
The Building Energy Lab runs this problem on a detached house with five appliances under a
time-of-use tariff, and compares the plan with an uncontrolled house and with a rule-based
battery. Two experiments belong to this chapter: removing the terminal value and watching the
battery end the day empty, and asking the house for a flexibility service --- a power cap or a
power curve --- with a fee for delivering it and a penalty for failing. Measured
demonstrations of predictive control on real buildings exist, on offices, with savings of the
order of tens of percent on the energy for heating and cooling.
\end{labbox}

\subsection{An aggregator of electric vehicles}
\label{sec:aggregator}

\paragraph{Context.} Electricity is traded on wholesale markets by generators, retailers and
the market operator. Next to them, the transmission and distribution system operators buy on
the \emph{ancillary services} or \emph{flexibility} markets what they need to keep frequency
and voltage at their references. The providers are large consumers or \emph{aggregators} of
many small resources, managed as one entity to reach the minimum bidding size --- of the order
of a megawatt for a given window. A fleet of charging vehicles is such an aggregation: the goal
of the aggregator is to make the \emph{total} charging power of the fleet track a profile agreed
with the market, while every driver still receives the charge requested.

\paragraph{Model.} For each vehicle $p$ connected during the window, a charging and a
discharging power, $P^{ch}_{\ik{i},p}, P^{dis}_{\ik{i},p} \ge 0$, with their indicators and
mutual exclusion; the stored energy with the efficiencies,
$x_{\ik{i+1},p} = x_{\ik{i},p} + \eta^{ch}\,T\,P^{ch}_{\ik{i},p} - \tfrac{1}{\eta^{dis}}\,T\,P^{dis}_{\ik{i},p}$;
and the fact that a charger cannot deliver an arbitrarily small power, the semi-continuous
constraint \eqref{eq:semicont}, $\delta^{ch}_{\ik{i},p}P^{ch,\min}_p \le P^{ch}_{\ik{i},p} \le
\delta^{ch}_{\ik{i},p}P^{\max}_p$. The total power is $P_{\ik{i}} = \sum_p \big(P^{ch}_{\ik{i},p}
- P^{dis}_{\ik{i},p}\big)$.

\paragraph{User preferences.} Each driver states a final charging time $F_p$ and a desired
state of charge $x^{ref}_p$. If the vehicle leaves inside the window, $F_p \le k + N$, the
request is a constraint, $x_{F_p|k,p} \ge x^{ref}_p$; if it leaves after the window, the request
enters the cost as a term $(x_{\ik{N},p} - x^{ref}_p)^2$. The constraint can be impossible to
meet --- a vehicle that arrives too empty or stays too little --- and it is therefore
\emph{relaxed} as in \cref{sec:soft}:
\begin{equation}
  x_{F_p|k,p} \ge x^{ref}_p - \varepsilon_p, \qquad \varepsilon_p \ge 0 .
\end{equation}

\paragraph{Objective.}
\begin{equation}
  J = d_1 \sum_{i=0}^{N-1}\big(P_{\ik{i}} - P^{ref}_{\ik{i}}\big)^2
    + d_2 \sum_{p:\,F_p > k+N}\big(x_{\ik{N},p} - x^{ref}_p\big)^2
    + d_3 \sum_{p} \varepsilon_p^2 .
\end{equation}
The weight $d_3$ decides the behaviour: small, and the controller disappoints a driver to follow
the grid better; large, and it disappoints a driver only when the request cannot be met at all.
Sweeping it traces the Pareto frontier between the two goals.

\begin{labbox}
The PEV Aggregator Lab runs the problem in closed loop over a day, with the request for a stop
in the aggregate power arriving at a chosen hour: until then the controller spreads the charge,
and at that hour the plan is rearranged. In Italy the role described here is regulated: Terna's
UVAM pilot admits aggregates of loads, storage and vehicles of at least one megawatt into the
balancing market, paid for availability and for the energy actually moved.
\end{labbox}

\subsection{Traffic lights at an intersection}
\label{sec:traffic}

\paragraph{Problem.} One intersection with a set $\mathcal{J}$ of traffic lights, one per
approaching road (\cref{fig:intersection}). Two goals belong to two different people: the
traffic manager wants short queues, the drivers want short waiting times.

\begin{figure}[t]
\centering
\begin{tikzpicture}[x=1cm,y=1cm,>=Stealth]
  \fill[gray!15] (-3,-0.5) rectangle (3,0.5);
  \fill[gray!15] (-0.5,-2.5) rectangle (0.5,2.5);
  \draw[gray!50,dashed] (-3,0) -- (-0.5,0);
  \draw[gray!50,dashed] (0.5,0) -- (3,0);
  \draw[gray!50,dashed] (0,-2.5) -- (0,-0.5);
  \draw[gray!50,dashed] (0,0.5) -- (0,2.5);
  \draw[thick] (-0.55,-0.5) -- (-0.55,0.5);
  \draw[thick] (-0.5,-0.55) -- (0.5,-0.55);
  \foreach \i in {0,...,3} \fill[mpc] (-0.95-0.55*\i,-0.17) rectangle (-0.6-0.55*\i,0.17);
  \foreach \i in {0,1} \fill[sap!70] (-0.17,-0.95-0.55*\i) rectangle (0.17,-0.6-0.55*\i);
  \node[font=\small,mpc] at (-2.1,0.75) {queue $m_{\ik{i},1}$};
  \node[font=\small,sap!80!black] at (1.35,-1.6) {queue $m_{\ik{i},2}$};
  \draw[->,thick] (2.4,0) -- (2.9,0);
  \draw[->,thick] (0,2.4) -- (0,2.9);
  \node[circle,draw,fill=ess,minimum size=5pt,inner sep=0] at (-0.85,0.85) {};
  \node[font=\scriptsize,above] at (-0.85,0.95) {light 1};
  \node[circle,draw,fill=sap,minimum size=5pt,inner sep=0] at (0.85,-0.85) {};
  \node[font=\scriptsize,right] at (0.95,-0.85) {light 2};
\end{tikzpicture}
\caption{The intersection of the fourth case study: two one-way roads, one light each, queues
counted at the stop lines.}
\label{fig:intersection}
\end{figure}

\paragraph{Control variables.} The colours: for each light $J$ three binaries
$g_{\ik{i},J}, y_{\ik{i},J}, r_{\ik{i},J} \in \{0,1\}$, exactly one of which is on,
\begin{equation}
  g_{\ik{i},J} + y_{\ik{i},J} + r_{\ik{i},J} = 1 .
\end{equation}
The admissible transitions are green $\to$ yellow $\to$ red $\to$ green, and each forbidden
transition is a linear inequality: green cannot be followed by red, red by yellow, yellow by
green,
\begin{equation}
  r_{\ik{i+1},J} \le 1 - g_{\ik{i},J}, \qquad
  y_{\ik{i+1},J} \le 1 - r_{\ik{i},J}, \qquad
  g_{\ik{i+1},J} \le 1 - y_{\ik{i},J}.
\end{equation}
The same inequalities, written between the first step of the window and the colour actually
displayed at $k-1$, tie the plan to the current state of the lights. Conflicting lights cannot
be green or yellow at the same time, $\sum_{J}\big(g_{\ik{i},J} + y_{\ik{i},J}\big) \le 1$.

\paragraph{The yellow.} Safety requires a minimum duration of the yellow. A timer counts how
long the yellow has been on, $t^y_{\ik{i+1},J} = (t^y_{\ik{i},J} + T)\,y_{\ik{i},J}$, a product
of a real variable and a binary linearized as in \cref{sec:product}, and the light cannot turn
red before the timer reaches the minimum:
\begin{equation}
  r_{\ik{i+1},J} \le \frac{t^y_{\ik{i},J}}{t^{y,\min}_J} + \big(1 - y_{\ik{i},J}\big).
\end{equation}
While the timer is short of the minimum the first term is strictly less than one, and since
$r$ is binary the constraint forces $r_{\ik{i+1},J} = 0$; the second term makes the constraint
irrelevant when the light is not yellow.

\paragraph{Queues and waiting times.} Two models. The \emph{aggregate} model keeps one state
per light, the number of vehicles waiting, with the dynamics
\begin{equation}
  m_{\ik{i+1},J} = m_{\ik{i},J} + v^{in}_{\ik{i},J}\,T - v^{out}_{J}\,T\,g_{\ik{i},J},
  \qquad m_{\ik{0},J} = m_{k,J} \text{ measured},\quad m \ge 0,
\end{equation}
where $v^{in}$ is the arrival rate and $v^{out}$ the discharge rate on green; the constraint
$m \ge 0$ is an approximation of the fact that a queue cannot become negative. Average waiting
times then follow from Little's law. The \emph{detailed} model keeps position, speed and
acceleration of every vehicle, the acceleration being a model of the driver for human-driven
cars and a command for automated ones, and a binary $\delta_{\ik{i},v}$ that says whether the
vehicle is still in the queue, so that its waiting time grows as $w_{\ik{i+1},v} = (w_{\ik{i},v}
+ T)\,\delta_{\ik{i},v}$ and the queue length is the sum of the indicators.

\paragraph{Objective.} On the queues, $J^m = \sum_i \sum_J m_{\ik{i},J}^2$: the square
shortens the queues and balances them, since for a given total it prefers two queues of five
to one of ten. On the waiting times the goal is to minimize the \emph{largest} of them, written
with an epigraph variable $\bar w_{\ik{i}}$, $w_{\ik{i},v} \le \bar w_{\ik{i}}$ for every vehicle
present, and $J^w = \sum_i \bar w_{\ik{i}}^2$ in the cost. The problem is a mixed-integer
quadratic program; a city is many intersections tied by coupling constraints, which is where
centralized and distributed MPC return.

\begin{labbox}
The Traffic Light Lab runs the intersection of \cref{fig:intersection} with the same vehicles
arriving on both sides of the screen, a fixed-time plan on the left and the MPC on the right,
with queues and waiting times compared as the day goes on. With two lights the transition
constraints leave few enough admissible colour sequences over the window for the page to
enumerate them and find the best one exactly; the mixed-integer formulation is what is needed
when the intersection is larger. Systems of this kind run in cities: Surtrac in Pittsburgh re-plans every intersection on
a rolling horizon about once a second, and UTOPIA, born in Turin, has been in operation since
the nineteen-eighties.
\end{labbox}

\subsection{Task scheduling in an assembly line}
\label{sec:scheduling}

\emph{This section covers what was done in class up to the lesson of 26 October; it will be
completed after the next lesson.}

\paragraph{Problem.} A \emph{job} --- the assembly of a vehicle --- is made of elementary
\emph{tasks}, executed on \emph{workstations} and requiring \emph{resources}: tools, people,
energy. Among the possible goals --- cost, quality, scrap, balanced use of the workstations ---
the one taken here is the total completion time of the job, the \emph{makespan}, read on a
Gantt chart. The orders of magnitude are what makes the problem hard: ten thousand tasks over
sixty days.

\paragraph{Sets and parameters.} The tasks $\mathcal{T}$, the workstations $\mathcal{W}$ and,
for each task, the subset $\mathcal{W}_t \subseteq \mathcal{W}$ of workstations that can execute
it. Each task $t$ carries an earliest start $S_t$, a deadline $F_t$, an estimate $d_t$ of the
number of steps it needs, and the number $R^{\#}_{t,r}$ of resources of each type $r \in
\mathcal{R}$ it consumes; $R^{a}_r$ is the number available.

\paragraph{Decision variables and dynamics.} One binary per task, workstation and time,
$e_{\ik{i},t,w} \in \{0,1\}$, equal to one when task $t$ runs at time $i|k$ on workstation $w$.
A task runs on at most one workstation at a time,
\begin{equation}
  \sum_{w \in \mathcal{W}_t} e_{\ik{i},t,w} \le 1, \qquad
  e_{\ik{i},t} = \sum_{w \in \mathcal{W}_t} e_{\ik{i},t,w} \in \{0,1\},
\end{equation}
where the auxiliary $e_{\ik{i},t}$ says whether the task is running. The state is the time left
to complete each task, decremented by one unit for every step actually worked,
\begin{equation}
  d_{\ik{i+1},t} = d_{\ik{i},t} - \sum_{w \in \mathcal{W}_t} e_{\ik{i},t,w},
\end{equation}
and the deadline is $d_{F_t,t} = 0$ for every task whose deadline falls inside the window.

\paragraph{Precedence constraints.} The relations that make a schedule a schedule:
finish-to-start, finish-to-finish, start-to-start, start-to-finish. The first, task $t_2$ cannot
run until $t_1$ is finished, is the one to be written as an inequality in the next lesson,
together with the resource constraints and the objective.

\begin{reading}
For the modelling constructions of all five case studies, \refkey{Hax\_1977}, chapter 9, and
\refkey{Borrelli\_2017}. For the market context of the aggregator, \refkey{Kirschen\_2019},
chapter 3 on markets for electrical energy and chapter 6 on power system operation; background
reading. For the traffic and the scheduling applications, the papers \refkey{Liberati\_2025a}
and \refkey{Liberati\_2025b} show the same formulations on real problems; background reading,
not examination material.
\end{reading}
